% Variante : plusieurs modules de macros mathématiques chargés à la carte.
% math={control,measure} — la virgule DOIT être protégée par des accolades :
% sans elles, \usepackage coupe sur la virgule et seule la première valeur
% atteint la clef (silencieusement). Le corps ci-dessous exerce quelques
% macros propres à chaque module ; base est toujours chargé.
\documentclass[11pt]{ocots-book}
\usepackage[lang=fr, theme=ocots, solutions=inline, math={control,measure}]{ocots}
\lstset{language=[LaTeX]TeX}

\title{Variante : modules mathématiques à la carte}
\author{Prénom \textsc{Nom}}
\date{\today}

\newcommand{\ocotsvariantnote}[1]{\begin{quote}\itshape Variante : #1\end{quote}}

\begin{document}

\begin{lstlisting}
\documentclass[11pt]{ocots-book}
\usepackage[lang=fr, theme=ocots, math={control,measure}]{ocots}
\end{lstlisting}
\ocotsvariantnote{\texttt{math=\{control,measure\}} (le défaut est
\texttt{math=} vide, c'est-à-dire \texttt{base} seul) : chaque valeur charge
\texttt{tex/math/ocots-math-<module>.sty}. Les accolades autour de la liste
sont obligatoires — la virgule non protégée est avalée par
\texttt{\textbackslash usepackage} avant d'atteindre l'option.}

\chapter{Module \texttt{control}}
\section{Solutions remarquables et homotopie}

Solutions remarquables barrées~: $\rsol$, $\fsol$, $\usol$, $\psol$, instant
final $\tfsol$. Paramètres d'homotopie~: $\sbar$, $\cbar$, $\lbar$, $\laz$,
application $\hom$, sphère $\sphere$.

Versions tildées (chemin d'homotopie)~: $\xt$, $\ut$, $\pt$, $\ft$, $\ct$,
$\Phit$, ensembles $\Et$, $\Ucalt$, $\Acalt$.

Crochet de dualité, application exponentielle, arc~:
\[
    \crochetDualite{p}{x}, \qquad \expmap{x}{t}{h}, \qquad \arc.
\]

Logiciels~: \hampath, \bocop, \cotcot, \nutopy, \controltoolbox, \tapenade,
\lapack, \minpack.

\chapter{Module \texttt{measure}}
\section{Tribus et fonctions mesurables}

Tribu borélienne $\Borel$, ensemble des parties $\PowerSet$, fonctions
mesurables $\Measurable$, fonctions étagées $\Simple$.

Classe d'équivalence $\eqclass{f}$ et convergence presque partout
$u_n \convae u$.

\end{document}
